\section{Introduction}

Gaussian elimination \index{Gaussian elimination} is one of the fundamental algorithms of linear algebra. Gauss developed it for solving systems of linear equations arising from geodetic surveying \index{geodesy} and least squares problems. The method systematically eliminates variables to reduce a system to triangular form.

In this chapter, we cover:
\begin{itemize}
    \item The forward elimination process
    \item Partial pivoting \index{pivoting!partial} for numerical stability
    \item LU factorization \index{LU factorization} and its computational complexity
    \item Back substitution for solving triangular systems
    \item Applications to least squares and linear systems
    \item Numerical considerations and error analysis
    \item Sparse matrix \index{sparse matrix} factorization and fill-in
    \item Connections to orthogonal factorizations (QR, Cholesky)
\end{itemize}

For $Ax = b$ with $A \in \R^{n\times n}$:
\begin{itemize}
    \item Forward elimination transforms $A$ to upper triangular form
    \begin{itemize}
        \item Row operations preserve the solution set
        \item Each step eliminates one variable below the pivot
    \end{itemize}
    \item Back substitution solves the triangular system $Ux = c$
    \begin{itemize}
        \item Proceeds from $x_n$ upward to $x_1$
        \item Cost is $O(n^2)$ operations
    \end{itemize}
\end{itemize}
\begin{itemize}
    \item LU factorization: $A = LU$ for efficient repeated solves
    \item Pivoting strategies control numerical growth
    \item Complexity: $\frac{2n^3}{3}$ flops for factorization
\end{itemize}

The standard Gaussian elimination algorithm performs:
\begin{itemize}
    \item For each column $k = 1, \ldots, n-1$:
    \item Eliminate entries below the pivot $A_{kk}$
    \item Update the remaining submatrix
\end{itemize}

The cost is dominated by the triangular array updates. For an $n \times n$ matrix, the total operation count is approximately $\frac{2n^3}{3}$ floating-point operations.

For $m$ right-hand sides, the LU factorization costs $\frac{2n^3}{3}$ flops once, then each solve costs $2n^2 - \frac{2n}{3}$ flops. For sparse matrices, specialized algorithms (e.g., sparsity-preserving permutations) can dramatically reduce the count, though fill-in may complicate the analysis.

\section{Forward Elimination}

The forward elimination phase transforms the system $Ax = b$ into an equivalent upper triangular system. At step $k$, we use row $k$ as the pivot row to eliminate all entries below $A_{kk}$ in column $k$.

The elementary row operation is:
\[
R_i \leftarrow R_i - \frac{A_{ik}}{A_{kk}} R_k, \quad i = k+1, \ldots, n
\]

After elimination, the matrix has zeros below the diagonal in column $k$. Repeating for all $k$ produces an upper triangular matrix $U$.

\begin{definition}[Multiplier \index{multiplier (Gaussian elimination)}]
The multiplier $m_{ik} = A_{ik}/A_{kk}$ is stored in $L$ for efficient LU factorization.
\end{definition}

\section{Partial Pivoting}

Gaussian elimination without pivoting can be numerically unstable. If a pivot element $A_{kk}$ is small relative to other entries in column $k$, the multipliers can become large, amplifying rounding errors.

Partial pivoting \index{pivoting!partial} addresses this by selecting the largest available entry in column $k$ (from rows $k$ through $n$) as the pivot. We swap row $k$ with the pivot row before elimination.

The partial pivoting strategy guarantees that all multipliers satisfy $|m_{ik}| \leq 1$, providing a growth factor bound of $\rho \leq 2^{n-1}$ in the worst case, though typical growth is much smaller.

\begin{definition}[Growth Factor \index{growth factor}]
The growth factor $\rho = \max_{i,j,k} |A_{ij}^{(k)}| / \max_{i,j} |A_{ij}|$ measures the amplification of matrix entries during elimination.
\end{definition}

Gaussian elimination with partial pivoting is \textbf{backward stable}: the computed solution $\hat{x}$ satisfies $(A + \delta A)\hat{x} = b + \delta b$ with $\|\delta A\|_\infty \leq \gamma_n \|A\|_\infty$ and $\|\delta b\|_\infty \leq \gamma_n \|b\|_\infty$, where $\gamma_n = \frac{nu}{1 - nu}$ and $u$ is the machine precision.

However, the \textbf{growth factor} $\rho = \max_{i,j,k} |A_{ij}^{(k)}| / \max_{i,j} |A_{ij}|$ can be as large as $2^{n-1}$ in pathological cases (Wilkinson, 1961). In practice, $\rho$ is typically $O(n^{1/3})$ for random matrices. The error bound is:
\[
\frac{\|x - \hat{x}\|}{\|x\|} \leq \gamma_n \cdot \kappa_\infty(A) \cdot \rho
\]
where $\kappa_\infty(A) = \|A\|_\infty \|A^{-1}\|_\infty$ is the condition number.

\textbf{Key insight}: Partial pivoting controls the growth factor for most practical matrices. Complete pivoting (searching both rows and columns) gives $\rho \leq n^{1/2}$ but at twice the overhead and no proven advantage in practice.

\section{LU Factorization}

The LU factorization \index{LU factorization} decomposes $A$ into the product of a lower triangular matrix $L$ and an upper triangular matrix $U$:
\[
A = LU
\]
where $L$ has ones on the diagonal (unit lower triangular) and $U$ is upper triangular.

With partial pivoting, we obtain the factorization:
\[
PA = LU
\]
where $P$ is a permutation matrix encoding the row swaps performed during elimination.

The LU factorization is computationally efficient for solving multiple right-hand sides:
\begin{itemize}
    \item Factorization: $\frac{2n^3}{3}$ flops (done once)
    \item Solve $Ly = Pb$: $n^2$ flops
    \item Solve $Ux = y$: $n^2$ flops
    \item Total for $m$ right-hand sides: $\frac{2n^3}{3} + 2mn^2$ flops
\end{itemize}

For symmetric positive definite matrices, the Cholesky factorization $A = LL^T$ is more efficient, requiring only $\frac{n^3}{3}$ flops and being inherently stable without pivoting.

For general matrices, the QR factorization $A = QR$ with orthogonal $Q$ provides superior numerical stability but at $O(n^3)$ cost with a larger constant.

\section{Back Substitution}

Once the system is in upper triangular form $Ux = c$, back substitution proceeds from the last equation upward:

\[
x_i = \frac{1}{U_{ii}}\left(c_i - \sum_{j=i+1}^{n} U_{ij}x_j\right), \quad i = n, n-1, \ldots, 1
\]

Each step requires $n-i$ multiplications and $n-i$ subtractions, for a total of $n^2 - n$ flops. The back substitution is numerically stable if the diagonal entries of $U$ are not too small relative to the matrix entries.

\section{Numerical Considerations}

Gaussian elimination has several important limitations:
\begin{itemize}
    \item \textbf{Pivoting necessity}: Without pivoting, small pivots can cause catastrophic cancellation
    \item \textbf{Fill-in}: For sparse matrices \index{sparse matrix}, elimination can create new nonzeros in previously zero positions
    \item \textbf{Conditioning}: Ill-conditioned systems amplify rounding errors regardless of pivoting
    \item \textbf{Complexity}: The $\frac{2n^3}{3}$ flop count limits direct methods to moderate $n$
    \item \textbf{Sparsity preservation}: Approximate minimum degree ordering reduces fill-in
\end{itemize}

In \textbf{practice}:
\begin{itemize}
    \item For dense systems, LAPACK's \texttt{dgetrf}/\texttt{dgetrs} are the standard implementation
    \item For sparse systems, approximate minimum degree ordering reduces fill-in
    \item For structured matrices (banded, block-diagonal), specialized algorithms exploit the structure
    \item Iterative methods (Conjugate Gradient, GMRES) become preferable for very large sparse systems
\end{itemize}

\section{Python Implementation}

The chapter's companion code in \texttt{python/gaussian\_elim.py} implements:
\begin{itemize}
    \item \texttt{gaussian\_elimination(A, b)}: Basic implementation with optional pivoting
    \item \texttt{lu\_decomposition(A)}: Returns $L$, $U$, and permutation vector
    \item \texttt{cholesky(A)}: Cholesky factorization for SPD matrices
    \item \texttt{back\_substitution(U, b)}: Solves upper triangular system
\end{itemize}

The implementation follows the mathematical description closely, with careful attention to numerical stability through partial pivoting.

\section{Exercises}

\begin{exercise}
Derive the operation count for Gaussian elimination: show that the total is $\frac{2n^3}{3} + O(n^2)$ flops.
\end{exercise}

\begin{exercise}
Implement Gaussian elimination with partial pivoting from scratch and verify against \texttt{numpy.linalg.solve}.
\end{exercise}

\begin{exercise}
Construct a $3 \times 3$ matrix where the growth factor $\rho = 4$.
\end{exercise}

\begin{exercise}
Show that for a symmetric positive definite matrix, Cholesky factorization \index{Cholesky factorization} requires $\frac{n^3}{3}$ flops.
\end{exercise}

\begin{exercise}
Implement a sparse Gaussian elimination with approximate minimum degree ordering.
\end{exercise}

\begin{exercise}
Prove that Gaussian elimination with partial pivoting is backward stable for all nonsingular matrices.
\end{exercise}

\section{Further Reading}

\begin{itemize}
    \item \cite{wilkinson} -- Wilkinson, \textit{Rounding Errors in Algebraic Processes}
    \item \cite{golub-vanloan} -- Golub \& Van Loan, \textit{Matrix Computations}
    \item \cite{trefethen-bau} -- Trefethen \& Bau, \textit{Numerical Linear Algebra}
\end{itemize}
